% ============================================================
%  Section 8: Conclusion
%  H&E to Sirius Red Virtual Staining Paper
% ============================================================

\section{Conclusion}
\label{sec:conclusion}

We conduct a full-factorial scaling study across six unsupervised I2I architectures,
three generator capacities, and three training data fractions, yielding 54
configurations.
Each configuration is evaluated on perceptual metrics (Patch-SSIM, LPIPS, FID) and a
task-specific CPA~MAE grounded in the paired tissue-block structure.
Building on this study, deep ensembles of ten models per best-performing family
configuration yield spatially resolved epistemic uncertainty maps, constituting the
first systematic comparison of epistemic uncertainty across unsupervised stain-to-stain
architectures.
To support this benchmark, we release the paired H\&E and Sirius Red WSI dataset from
the mouse bile-duct ligation experiment of~\cite{ghallab2025asbt}.

The scaling study reveals that perceptual metrics saturate within a narrow band across
the majority of well-behaved configurations and correlate only moderately with CPA~MAE
($r_s{=}0.55$--$0.64$), confirming that task-specific evaluation is indispensable for
model selection.
CPA~MAE ranges from $0.008$ to $0.171$ across the 54 configurations (a 21-fold span),
with medium-capacity generators trained on the full dataset performing most reliably
across families.
MUNIT is the most stable architecture, maintaining a narrow CPA~MAE range of
$0.040$--$0.059$ across all nine scaling conditions.
DCLGAN's large generator collapses at reduced data fractions, while UVCGAN's
medium and large variants fail across all data fractions; both underscore that
over-parameterisation can degrade translation as severely as under-training.
CycleDiffusion produces competitive CPA~MAE despite systematically lower Patch-SSIM
than GAN-based families, demonstrating that perceptual sharpness and biological
accuracy are not equivalent objectives.
Per-family ensemble variance separates the six families along an axis distinct from
CPA~MAE: DCLGAN pairs low task error with the narrowest inter-seed spread
(IQR $0.043$) and the most informative variance signal among GAN families, while
CycleDiffusion records the lowest mean ensemble variance ($3.16$) but the widest spread
and is uncorrelated with cycle-reconstruction error.
Uncertainty signals are informative only at tile resolution; pixel-level uncertainty
does not reliably co-locate with high-error regions within a tile
(Supplementary Sec.~\ref{sec:supp-unc-results}).
These findings indicate that perceptual quality, task accuracy, and epistemic
uncertainty are three independent evaluation axes for unsupervised stain translation.

The dataset, model implementations, and evaluation code are released publicly to
support reproducible benchmarking and future work on virtual histochemical staining.
